\documentclass[11pt]{article}
\usepackage{amsmath,amssymb,amsthm,geometry,booktabs,hyperref}
\geometry{margin=2.5cm}
\newtheorem{theorem}{Theorem}
\newtheorem{lemma}{Lemma}
\newtheorem{proposition}{Proposition}
\newtheorem{corollary}{Corollary}
\theoremstyle{definition}\newtheorem{definition}{Definition}
\theoremstyle{remark}\newtheorem{remark}{Remark}
\newcommand{\Q}{\mathcal Q}
\newcommand{\R}{\mathbb R}
\newcommand{\Pg}{P_{\rm g}}
\title{Optimal Bell witnesses under detection loss lie in the normal cone of their own data}
\author{Draft note for the NDWU project (lossycone)}
\date{23 October 2026}
\begin{document}
\maketitle

\noindent\textbf{What this note proves.} Theorem~\ref{thm:own} is a general and elementary statement: the Bell expression that certifies the most randomness from the best honest device is maximised, over all honest devices, by that device's own behaviour. In geometric language it lies in the normal cone of the observed behaviour. Its converse is a practical test: if some honest device scores higher on your best expression, switching to that device certifies strictly more randomness with the same expression. Proposition~\ref{prop:reduce} reduces the optimal devices to two-qubit strategies and identifies the relevant cone with the ideal normal cone of the paper, transported through the loss map. Theorems~\ref{thm:cones} and~\ref{thm:rates} are computer-assisted: at the optimal devices we found for six efficiencies the ideal normal cone equals the explicit polygon $\mathcal C(q)$, and a Bell expression taken \emph{exactly} from that cone certifies $9$--$18\%$ more min-entropy per round than every CHSH protocol. Theorem~\ref{thm:cover} extends the rate statement to every efficiency in $[0.75,0.98]$. The Audit section lists what is pen-and-paper, what is exact or interval computation, and what is numerical.

\section{Setting}
Behaviours are $z=(\alpha_0,\alpha_1,\beta_0,\beta_1,E_{00},E_{01},E_{10},E_{11})\in\R^8$ and Bell expressions $g\cdot z$, as in the paper. $\Q$ is the quantum set, $h_A(g)=\max_{z\in A}g\cdot z$, and $N_A(p)=\{g:\ g\cdot p=h_A(g)\}$ is the normal cone of a convex set $A$ at $p\in A$. $h_L(g)$ is the maximum over the 16 deterministic points.

\paragraph{Detection loss.} Each detector clicks with probability $\eta\in(0,1]$ and a missed click is recorded as $+1$. The recorded behaviour is $L_\eta(z)=M_\eta z+m_\eta$ with
\[
\alpha_x'=\eta\alpha_x+1-\eta,\quad \beta_y'=\eta\beta_y+1-\eta,\quad E'_{xy}=\eta^2E_{xy}+\eta(1-\eta)(\alpha_x+\beta_y)+(1-\eta)^2 .
\]
$M_\eta$ is triangular with diagonal $(\eta,\eta,\eta,\eta,\eta^2,\eta^2,\eta^2,\eta^2)$, so $L_\eta$ is an invertible affine map. The honest behaviours at efficiency $\eta$ form the compact convex set $\Q_\eta=L_\eta(\Q)\subseteq\Q$. Two facts are used below. First, $L_a\circ L_b=L_{ab}$ (a click survives both stages with probability $ab$, and $+1$ stays $+1$). Second, for $\eta\le\eta'$, $\Q_\eta=L_{\eta'}(L_{\eta/\eta'}(\Q))\subseteq\Q_{\eta'}$: a lab with efficiency $\eta'$ can discard each click with probability $1-\eta/\eta'$ and so reproduce every honest behaviour at efficiency $\eta$.

\paragraph{Randomness.} Eve guesses Alice's outcome $A_1$. Conditioning on Eve's guess $e$ splits any attack into quantum behaviours $z_e\in\Q$ with weights $t_e\ge0$, $\sum_et_e=1$, and her success is $\sum_et_e\Pi_e(z_e)$ with $\Pi_e(z)=\frac12(1+e\alpha_1)$. The \emph{full-statistics} guessing probability of an observed behaviour $p$ is
\[
\Pg(p)=\sup\Bigl\{\textstyle\sum_et_e\Pi_e(z_e):\ \sum_et_ez_e=p\Bigr\},
\]
and the \emph{single-witness} guessing probability for a Bell expression $g$ and an observed score $\omega$ is
\[
\Pg(g,\omega)=\sup\Bigl\{\textstyle\sum_et_e\Pi_e(z_e):\ \sum_et_e\,g\cdot z_e\ge\omega\Bigr\}.
\]
$H_{\min}=-\log_2\Pg$. Everything below holds verbatim if $\Q$ in these two definitions is replaced by any closed convex set containing $\Q$, such as an NPA level; the honest set $\Q_\eta$ is not changed. A witness $g$ is \emph{tight at $p$} if $\Pg(g,g\cdot p)=\Pg(p)$: the single Bell value $g\cdot p$ certifies as much as the whole behaviour.

\section{General results (pen and paper)}
\begin{lemma}[Elementary properties]\label{lem:props}
(a) $\Pg(g,\omega)=\sup\{\Pg(q):\ q\in\Q,\ g\cdot q\ge\omega\}$; in particular $\Pg(p)\le\Pg(g,g\cdot p)$ for every $g$.
(b) $\Pg$ is concave on $\Q$, and $\omega\mapsto\Pg(g,\omega)$ is concave and non-increasing on $(-\infty,h_\Q(g)]$.
(c) $\Pg(g,\omega)=1$ for $\omega\le h_L(g)$.
(d) If $\Pg(g,\omega^*)<1$, then $\omega^*>h_L(g)$ and, with $\kappa=(1-\Pg(g,\omega^*))/(\omega^*-h_L(g))>0$,
\[
\Pg(g,\omega)\le\Pg(g,\omega^*)-\kappa(\omega-\omega^*)\qquad\text{for all }\omega\ge\omega^*.
\]
In particular $\Pg(g,\cdot)$ is strictly decreasing on $[\omega^*,h_\Q(g)]$.
\end{lemma}
\begin{proof}
(a) A decomposition with $\sum_et_eg\cdot z_e\ge\omega$ is a decomposition of $q=\sum_et_ez_e\in\Q$, and $g\cdot q\ge\omega$; conversely every decomposition of such a $q$ is admissible for $\Pg(g,\omega)$. The second statement takes $q=p$.
(b) Mix two admissible decompositions $\{t_e,z_e\}$ and $\{t'_e,z'_e\}$ with weights $\lambda,1-\lambda$: the weights $\lambda t_e+(1-\lambda)t'_e$ and the behaviours $(\lambda t_ez_e+(1-\lambda)t'_ez'_e)/(\lambda t_e+(1-\lambda)t'_e)\in\Q$ form a decomposition. Since $t_e\Pi_e(z_e)$ and $t_e\,g\cdot z_e$ are linear in the pair $(t_e,t_ez_e)$, the objective and the score mix linearly. This gives concavity of $\Pg$ in $p$ and of $\Pg(g,\cdot)$ in $\omega$. Monotonicity is clear.
(c) A deterministic point $L$ with $g\cdot L=h_L(g)\ge\omega$ is admissible and Eve guesses it perfectly.
(d) By (c), $\omega^*>h_L(g)$. For $\omega>\omega^*$ write $\omega^*=\lambda h_L(g)+(1-\lambda)\omega$ with $\lambda\in(0,1)$. Concavity gives $\Pg(g,\omega^*)\ge\lambda+(1-\lambda)\Pg(g,\omega)$, which rearranges to the stated linear bound (the chord from $(h_L(g),1)$ through $(\omega^*,\Pg(g,\omega^*))$ lies above the graph beyond $\omega^*$).
\end{proof}

\begin{lemma}[Tight witnesses exist in the interior]\label{lem:exist}
If $p$ is an interior point of $\Q$, some $g$ is tight at $p$.
\end{lemma}
\begin{proof}
$\Pg$ is concave and finite on $\Q$ (Lemma~\ref{lem:props}(b)), so at an interior point it has a supergradient $s$: $\Pg(q)\le\Pg(p)+s\cdot(q-p)$ for all $q\in\Q$. Put $g=-s$. If $g\cdot q\ge g\cdot p$, then $\Pg(q)\le\Pg(p)$; by Lemma~\ref{lem:props}(a), $\Pg(g,g\cdot p)\le\Pg(p)$, and the reverse inequality is (a) again.
\end{proof}

\begin{theorem}[The optimal witness lies in the normal cone of its own data]\label{thm:own}
Let $K\subseteq\Q$ be a convex set of honest behaviours (for example $K=\Q_\eta$). Let $p^*\in K$ be a local minimiser of $\Pg$ on $K$ with $\Pg(p^*)<1$ (that is, the device certifies a positive amount of randomness and no nearby honest device certifies more), and let $g$ be tight at $p^*$. Then
\[
g\cdot p\le g\cdot p^*\quad\text{for all }p\in K,\qquad\text{i.e.}\qquad g\in N_K(p^*).
\]
\end{theorem}
\begin{proof}
Let $\omega^*=g\cdot p^*$. Then $\Pg(g,\omega^*)=\Pg(p^*)<1$. Suppose some $p\in K$ has $g\cdot p>\omega^*$. For $t\in(0,1]$, $p_t=p^*+t(p-p^*)\in K$ and $g\cdot p_t>\omega^*$, so by Lemma~\ref{lem:props}(a,d)
\[
\Pg(p_t)\le\Pg(g,g\cdot p_t)<\Pg(g,\omega^*)=\Pg(p^*).
\]
Letting $t\to0$ contradicts local minimality.
\end{proof}
\noindent No monotonicity assumption is needed: strict decrease of the score-to-guessing curve follows from concavity once the device certifies anything at all. The same proof applies to any certified entropy $R$ that is convex in the behaviour, vanishes at deterministic points, and is evaluated for a single score as $R(g,\omega)=\inf\{R(q):q\in\Q,\ g\cdot q\ge\omega\}$; this covers the von Neumann entropy $H(A_1|X{=}1,E)$, which is convex because Eve may keep a classical record of a mixture.

\begin{corollary}[Improvement test]\label{cor:improve}
Let $p\in K$ with $\Pg(p)<1$ and let $g$ be tight at $p$. If some $p'\in K$ has $g\cdot p'>g\cdot p$, then $p'$ certifies strictly more min-entropy than $p$, and it does so with the same Bell expression $g$: $\Pg(g,g\cdot p')<\Pg(g,g\cdot p)=\Pg(p)$.
\end{corollary}
\noindent In words: if your best Bell expression is not maximised by your own device, a better device exists and the same expression shows it.

\begin{proposition}[Quantitative version]\label{prop:quant}
Let $p^*$ minimise $\Pg$ on $K$ with $\Pg(p^*)<1$, and let $g$ be $\varepsilon$-tight at $p^*$, i.e.\ $\Pg(g,g\cdot p^*)\le\Pg(p^*)+\varepsilon$. Then $g\cdot p\le g\cdot p^*+\varepsilon/\kappa$ for all $p\in K$, with $\kappa$ as in Lemma~\ref{lem:props}(d) for $\omega^*=g\cdot p^*$.
\end{proposition}
\begin{proof}
For $p\in K$ with $g\cdot p>\omega^*$: $\Pg(p^*)\le\Pg(p)\le\Pg(g,g\cdot p)\le\Pg(g,\omega^*)-\kappa(g\cdot p-\omega^*)\le\Pg(p^*)+\varepsilon-\kappa(g\cdot p-\omega^*)$.
\end{proof}

\begin{proposition}[Reduction to two-qubit points and to the ideal cone]\label{prop:reduce}
Let $0<\eta\le1$.
(i) $\inf_{\Q_\eta}\Pg=\inf\{\Pg(L_\eta q):\ q \text{ an extreme point of }\Q\}$. If the infimum is attained and is $<1$, it is attained at $L_\eta(q^*)$ with $q^*$ a nonlocal extreme point of $\Q$; such $q^*$ is a qubit point in the sense of the paper (pure entangled two-qubit state, coplanar projective measurements, Bob's overlap $c_0\in(-1,1)$, all four slice constraints active).
(ii) For every $q\in\Q$, $N_{\Q_\eta}(L_\eta q)=M_\eta^{-\top}N_\Q(q)$.
(iii) Consequently, at such an optimal device every tight witness $g$ satisfies $M_\eta^\top g\in N_\Q(q^*)\subseteq\mathcal C(q^*)$: its ideal pull-back rewards the steering pattern of $q^*$ (slice cone), is stationary under Bob's overlap, and is not beaten by any deterministic strategy.
\end{proposition}
\begin{proof}
(i) $L_\eta$ is an affine bijection, so the extreme points of $\Q_\eta$ are the images of those of $\Q$. Every point of the compact convex set $\Q_\eta\subset\R^8$ is a convex combination $\sum_i\lambda_ie_i$ of finitely many extreme points (Minkowski), and concavity gives $\Pg(\sum\lambda_ie_i)\ge\sum_i\lambda_i\Pg(e_i)\ge\min_i\Pg(e_i)$. If $p^*$ attains the infimum, every $e_i$ with $\lambda_i>0$ attains it too. Extreme points of $\Q$ that are local are deterministic, and there $\Pg=1$; the nonlocal ones are qubit points by Lemma~``Scope'' of the paper.
(ii) $g\cdot L_\eta(z)=(M_\eta^\top g)\cdot z+g\cdot m_\eta$, so $g$ is maximised over $\Q_\eta$ at $L_\eta q$ iff $M_\eta^\top g$ is maximised over $\Q$ at $q$.
(iii) Theorem~\ref{thm:own} with $K=\Q_\eta$, (ii), and $N_\Q(q)\subseteq\mathcal C(q)$ (the three necessary conditions, Lemma ``Scope'' of the paper).
\end{proof}
\noindent So, without any computation, the search for the optimal witness under loss reduces to two-qubit devices and to the explicit polygon $\mathcal C(q)$ transported by $M_\eta^{-\top}$. The computer-assisted results below show that at our optimal devices nothing is lost: the polygon \emph{is} the normal cone.

\begin{remark}[When no tight witness exists]
If $p^*$ lies on the boundary of $\Q$, a tight witness need not exist. This is what happens at our level-3 numerical optima for $\eta=0.90$ and $0.95$: the ideal device there has $P(-1,-1|1,1)=0$ (a Hardy-type zero), so the observed behaviour lies on a positivity facet of $\Q$. The full-statistics duals then blow up along the facet direction, which is why those two certificates were ill-conditioned. Proposition~\ref{prop:quant} still applies to any $\varepsilon$-tight witness.
\end{remark}

\section{Computer-assisted results}
\paragraph{Devices.} For each $\eta\in\{0.75,0.80,0.85,0.90,0.95,0.98\}$ the device is $\cos\vartheta|00\rangle+\sin\vartheta|11\rangle$ with $A_x=\cos a_x Z+\sin a_x X$, $B_y=\cos b_y Z+\sin b_y X$, chosen numerically to maximise the min-entropy certified by the best expression in $M_\eta^{-\top}\mathcal C(q)$ (Nelder--Mead on NPA level 2 and 3 objectives). At $\eta=0.95$ the level-3 optimiser drifts towards a Hardy-type device (see the Remark), for which the cone is still verified but the NPA lower bound is far from the attack ($0.047$ vs $0.389$ bits); we therefore use the device obtained by warm-starting from the certified neighbouring grid point $\eta=0.949$, which is interior and certifies $0.3856$ bits. Every cosine and sine (of $2\vartheta$, $a_x$, $b_y$) is then replaced by a rational point on the unit circle, $(\frac{1-u^2}{1+u^2},\frac{2u}{1+u^2})$ with rational $u$. The ideal behaviour $q_\eta$, Bob's overlap $c_0$, the slice gradients, the overlap derivative and the polygon are then exactly rational. The parameters are in \texttt{data/t4/<eta>.json}.

\begin{theorem}[Exact normal cones at the optimal devices; computer-assisted]\label{thm:cones}
For each of the six devices $q_\eta$, all four slice constraints are active, and $N_\Q(q_\eta)=\mathcal C(q_\eta)$. Its normalised section is a polygon with $n_\eta$ rational vertices (Table~\ref{tab:main}), every vertex has all four slice multipliers $\mu_{a|x}>0$ and is tight on exactly two deterministic points. Hence $N_{\Q_\eta}(L_\eta q_\eta)=M_\eta^{-\top}\mathcal C(q_\eta)$ exactly.
\end{theorem}
\begin{proof}
``$\subseteq$'' is Lemma ``Scope'' of the paper. ``$\supseteq$'': the polygon is the section $\{\mu\ge0,\ \sum\mu=1,\ \mu\cdot m=0,\ (\sum\mu\gamma)\cdot L\le1\}$; its vertices are found by solving, in exact rationals, the $4\times4$ system for every pair of constraints and keeping the feasible solutions, so the vertex list is complete. For each vertex $V$ we give a global certificate $\delta(c)=d_0+d_1(c-c_0)+(c-c_0)^2[\sum_jR_jT_j(c)+\alpha+\beta c]$ with rational $R_j$ and the exact forced data $d_0,d_1$ (the unique slice dual at $c_0$ and the first-order equalities; the third first-order equation is implied exactly, which is overlap stationarity), and with an exact pin at $c=\pm1$ where both tight deterministic points have the same $B_0B_1$. Then $\Phi_c(\delta(c))\le1=V\cdot q_\eta$ for all $c\in[-1,1]$ is checked as follows: each $\|u(c)\|_{G_c}^2$ is an exact rational polynomial; $F=\partial_cF=0$ at $c_0$ and the endpoint values at $c=\pm1$ are verified in exact rational arithmetic; and a bisection of $[-1,1]$ in ball arithmetic (Arb, 160 bits, exact-midpoint Taylor shifts) accepts every interval by a direct enclosure, a mean-value bound, $\partial_c^2F>0$ on an interval containing $c_0$, or monotone approach to an endpoint. By the NDWU dual theorem and the generation theorem, $h_\Q(V)=V\cdot q_\eta$; convexity gives the whole cone. (ii) of Proposition~\ref{prop:reduce} transports the result.
\end{proof}

\begin{theorem}[Certified rates from exact cone witnesses; computer-assisted]\label{thm:rates}
For each $\eta$ in Table~\ref{tab:main}, the Bell expression $g_\eta=M_\eta^{-\top}\sum_i\lambda_iV_i$ (rational convex weights $\lambda_i$ over the vertices $V_i$ of Theorem~\ref{thm:cones}) lies exactly in $N_{\Q_\eta}(p_\eta)$, $p_\eta=L_\eta q_\eta$, and certifies the listed $H_{\min}(A_1|X{=}1,E)$ for every quantum device, in any dimension, that reaches its score; the interval is two-sided (an explicit attack attains the upper end). Every CHSH protocol at the same efficiency certifies at most the listed CHSH bound.
\end{theorem}
\begin{proof}
Lemma ``Two-sided bounds'' of the paper with exact-rational NPA sum-of-squares certificates at the exact score $\omega=g_\eta\cdot p_\eta$ (a rational number), and an explicit mixture of two-qubit strategies checked in ball arithmetic. The CHSH bound is the tight closed form at $S_{\rm up}=c_0+\beta$, with $\beta$ an exact-rational NPA certificate for the loss-transformed CHSH expression.
\end{proof}

\begin{table}[h]\centering
\begin{tabular}{cccccc}
$\eta$ & $n_\eta$ & $c_0$ & $H_{\min}$ (cone witness) & CHSH bound & gain\\\hline
0.75 & 5 & 0.4578 & $[0.009845,\,0.009847]$ & 0.009013 & $+9.2\%$\\
0.80 & 5 & 0.2826 & $[0.037859,\,0.037864]$ & 0.033445 & $+13.2\%$\\
0.85 & 7 & 0.1735 & $[0.095104,\,0.095105]$ & 0.081930 & $+16.0\%$\\
0.90 & 7 & 0.0935 & $[0.198626,\,0.198627]$ & 0.168313 & $+18.0\%$\\
0.95 & 6 & 0.1089 & $[0.385565,\,0.385566]$ & 0.330204 & $+16.7\%$\\
0.98 & 8 & 0.3556 & $[0.597367,\,0.597368]$ & 0.525646 & $+13.6\%$\\
\end{tabular}
\caption{Theorems~\ref{thm:cones} and~\ref{thm:rates}. $H_{\min}(A_1|X{=}1,E)$ in bits per round, lower ends rounded down and upper ends up; gain $=H_{\rm lo}/\bar H_{\rm CHSH}-1$, rounded down.}\label{tab:main}
\end{table}

\begin{theorem}[Every efficiency in an interval; computer-assisted]\label{thm:cover}
Let $\eta_1<\dots<\eta_N$ be the $N=110$ efficiencies of \texttt{grid.tsv}, from $\eta_1=0.750$ to $\eta_N=0.980$ (spacing $0.001$--$0.003$). At every $\eta_i$ the conclusions of Theorems~\ref{thm:cones} and~\ref{thm:rates} hold for the device stored in \texttt{data/t4/<eta>.json} (643 polygon vertices verified in total). Consequently, for \emph{every} $\eta\in[0.75,0.98]$ there is a two-qubit protocol at efficiency $\eta$ (the grid protocol at the nearest $\eta_i\le\eta$, with clicks discarded at random) which certifies more min-entropy per round than every CHSH protocol at the same efficiency, by at least
\[
5.3\%\ (\eta\in[0.75,0.80]),\quad 8.5\%\ ([0.80,0.85]),\quad 10.7\%\ ([0.85,0.90]),\quad 12.2\%\ ([0.90,0.95]),\quad 8.0\%\ ([0.95,0.98]).
\]
At the grid points the Bell expression lies exactly in the normal cone of the honest data; between grid points it need not (by Corollary~\ref{cor:improve} the emulated protocol is then not optimal, which only makes the bound conservative). At the grid points themselves the gain lies between $8.8\%$ and $17.6\%$.
\end{theorem}
\begin{proof}
Let $\eta_1<\dots<\eta_N$ be the certified grid. At each $\eta_i$ we have a device $p_i\in\Q_{\eta_i}$ and a witness certifying $H_i$ (as in Theorem~\ref{thm:rates}, with the exact cone verified as in Theorem~\ref{thm:cones}), and a certified CHSH bound $\bar H_i$. For $\eta\in[\eta_i,\eta_{i+1}]$: since $\Q_{\eta_i}\subseteq\Q_\eta$, a lab at efficiency $\eta$ can realise $p_i$ by discarding clicks, so $H_i$ is certifiable at $\eta$; and since $\Q_\eta\subseteq\Q_{\eta_{i+1}}$, every CHSH value at $\eta$ is at most $S_{\rm up}(\eta_{i+1})$, so every CHSH protocol certifies at most $\bar H_{i+1}$. Hence the gain on $[\eta_i,\eta_{i+1}]$ is at least $H_i/\bar H_{i+1}-1$.
\end{proof}

\section{Audit}
\paragraph{Pen and paper.} Lemma~\ref{lem:props}, Lemma~\ref{lem:exist}, Theorem~\ref{thm:own}, Corollary~\ref{cor:improve}, Proposition~\ref{prop:quant}, Proposition~\ref{prop:reduce}; $L_aL_b=L_{ab}$ and $\Q_\eta\subseteq\Q_{\eta'}$; the interval argument in Theorem~\ref{thm:cover}. Inputs from the paper: the generation theorem of Ref.~[NDWU] (used, not re-derived, exactly as in the paper), the NDWU dual theorem, Lemma ``Scope'' ($N_\Q\subseteq\mathcal C$), Lemma ``Two-sided bounds'', and the tightness of the CHSH closed form.
\paragraph{Exact rational computation.} Device data (rational unit-circle points), $q_\eta$, $c_0$, activity of the four slice constraints, $\gamma_{a|x}$, $\dot z$, $m$, the complete vertex enumeration of each polygon, $\mu>0$, the two tight points of each vertex, the forced certificate data $d_0,d_1$ and the consistency of the third first-order equation, $F(c_0)=F'(c_0)=0$, the endpoint values $F(\pm1)\ge0$ and the pins; the push $M_\eta^{-\top}$ (checked by $M_\eta^\top g=V$); the witness and its score $\omega$; all NPA sum-of-squares certificates (rational $LDL^\top$).
\paragraph{Interval (ball) arithmetic.} The bisection over $c\in[-1,1]$ for each vertex; the attacks (upper ends of the $H_{\min}$ intervals); $-\log_2$ of rational bounds and the CHSH closed form.
\paragraph{Numerical only.} (1) The device optimisation. Our devices are numerical optima, so the hypothesis of Theorem~\ref{thm:own} (local optimality) is \emph{not} verified for them, and the certified values are lower bounds on the best achievable rate, not the optimum. (2) Theorems~\ref{thm:cones}--\ref{thm:rates} are pointwise: they hold at the stated rational devices, not on an open set of devices; Theorem~\ref{thm:cover} covers every efficiency in the interval through the monotonicity argument, but the cone statement there is again at the grid devices. (3) The fact that the best expression in $M_\eta^{-\top}\mathcal C(q)$ reproduces the numerical full-statistics optimum (to $10^{-6}$ or better at five of six efficiencies) is a numerical observation consistent with Theorem~\ref{thm:own}. Nothing in the theorem statements depends on (1)--(3).

\end{document}
